% --- Archon physics formalization source begin ---
% archon:physics
% archon:covers PhyXMiniProblems/problem_phyx_mini_0173.lean
% archon:source-report reports/phyx_mini/problem_phyx_mini_0173.source.json

\chapter{Physics problem phyx\_mini\_0173}
\label{ch:PhyXMiniProblems_problem_phyx_mini_0173}

\paragraph{Problem source.}
\#\# Physical scenario

Two identical loudspeakers are located at points \textbackslash{}( A \textbackslash{}) and \textbackslash{}( B \textbackslash{}), \textbackslash{}( 2.00\textbackslash{},\textbackslash{}text\{m\} \textbackslash{}) apart. The loudspeakers are driven by the same amplifier and produce sound waves with a frequency of \textbackslash{}( 784\textbackslash{},\textbackslash{}text\{Hz\} \textbackslash{}). Take the speed of sound in air to be \textbackslash{}( 344\textbackslash{},\textbackslash{}text\{m/s\} \textbackslash{}). A small microphone is moved out from point \textbackslash{}( B \textbackslash{}) along a line perpendicular to the line connecting \textbackslash{}( A \textbackslash{}) and \textbackslash{}( B \textbackslash{}) (line \textbackslash{}( BC \textbackslash{}) in \textbackslash{}textbf\{figure\}).

\#\# Question

If the frequency is made low enough, there will be no positions along the line \textbackslash{}( BC \textbackslash{}) at which destructive interference occurs. How low must the frequency be for this to be the case?

\#\# Answer choices

- A: A: 86Hz

- B: B: 38.0Hz

- C: C: 89.6Hz

- D: D: 38.5Hz

\#\# Figure caption (auxiliary; use the image as primary evidence)

The image depicts a diagram involving two speakers, labeled A and B, and a microphone. 

- **Speaker A**: Positioned vertically above the horizontal line, marked at a height of 2.00 meters from Speaker B.
- **Speaker B**: Located on the horizontal line at point B.
- **Microphone**: Positioned along the horizontal line, indicated by a red icon. It is placed to the right of Speaker B at a distance labeled as "x."
- **Coordinate Points**: Three points are identified - A (associated with Speaker A), B (associated with Speaker B), and C (further along the horizontal path of the microphone).

The diagram shows a simple geometric setup, likely representing a situation involving sound waves or measurements related to acoustics, given the presence of speakers and a microphone.

\#\# Dataset metadata

Resonance and Harmonics; Physical Model Grounding Reasoning, Spatial Relation Reasoning

\paragraph{Recorded answer/context.}
A: A: 86Hz

\paragraph{Figure/image path.}
/root/proposal\_for\_physic/science-mango/phyx\_mini\_run/phyx\_data/test\_image/173.png

\paragraph{Formalization target.}
create a compiling Lean file with sorry bodies at `PhyXMiniProblems/problem\_phyx\_mini\_0173.lean`. The Lean declarations must preserve the physical quantities, dimensions or dimensional roles, figure labels, governing-law hypotheses, and final relation expressed by this problem.
Use Mathlib/Physlib names found through LeanExplore where available. If a physics API is missing, introduce faithful local abstractions rather than scalar placeholder aliases.

\begin{theorem}[Physics formalization target]
\label{thm:physics:phyx_mini_0173:target}
\lean{PhyXMiniProblems.ProblemPhyXMini0173.problem_phyx_mini_0173}
\uses{def:physics:phyx-mini-0173:phyxminiproblems-problemphyxmini0173-recordeddatasetanswer, def:physics:phyx-mini-0173:phyxminiproblems-problemphyxmini0173-twospeakerinterferencesetup, def:physics:phyx-mini-0173:phyxminiproblems-problemphyxmini0173-coherentlydrivenidenticalspeakers, def:physics:phyx-mini-0173:phyxminiproblems-problemphyxmini0173-hasdepictedgeometry, def:physics:phyx-mini-0173:phyxminiproblems-problemphyxmini0173-hasstatedreadouts, def:physics:phyx-mini-0173:phyxminiproblems-problemphyxmini0173-hasphysicalparameters, def:physics:phyx-mini-0173:phyxminiproblems-problemphyxmini0173-satisfiestwosourceacousticlaws, def:physics:phyx-mini-0173:phyxminiproblems-problemphyxmini0173-isnodestructiveinterferencecutoff}
The assigned autoformalize agent should translate this physics problem into Lean declarations in the covered file, with theorem and lemma proof bodies written as `by sorry`.
\end{theorem}
\begin{proof}
This is an autoformalization task, not a proof task. Produce statements that can later be proved by the physics prover without weakening the physical model.
\end{proof}

% --- Archon declaration topology begin ---
\section{Declaration topology}
\begin{definition}[Length Quantity]
  \label{def:physics:phyx-mini-0173:phyxminiproblems-problemphyxmini0173-lengthquantity}
  \lean{PhyXMiniProblems.ProblemPhyXMini0173.LengthQuantity}
  A signed physical length represented coherently in every unit system
\end{definition}
\begin{proof}
  This is the declared model definition of Length Quantity.
\end{proof}
\begin{definition}[Frequency Quantity]
  \label{def:physics:phyx-mini-0173:phyxminiproblems-problemphyxmini0173-frequencyquantity}
  \lean{PhyXMiniProblems.ProblemPhyXMini0173.FrequencyQuantity}
  A physical frequency, carrying the inverse-time dimension
\end{definition}
\begin{proof}
  This is the declared model definition of Frequency Quantity.
\end{proof}
\begin{definition}[meters Value]
  \label{def:physics:phyx-mini-0173:phyxminiproblems-problemphyxmini0173-metersvalue}
  \lean{PhyXMiniProblems.ProblemPhyXMini0173.metersValue}
  \uses{def:physics:phyx-mini-0173:phyxminiproblems-problemphyxmini0173-lengthquantity}
  Read a physical length in metres
\end{definition}
\begin{proof}
  This is the declared model definition of meters Value.
\end{proof}
\begin{definition}[hertz Value]
  \label{def:physics:phyx-mini-0173:phyxminiproblems-problemphyxmini0173-hertzvalue}
  \lean{PhyXMiniProblems.ProblemPhyXMini0173.hertzValue}
  \uses{def:physics:phyx-mini-0173:phyxminiproblems-problemphyxmini0173-frequencyquantity}
  Read a physical frequency in hertz
\end{definition}
\begin{proof}
  This is the declared model definition of hertz Value.
\end{proof}
\begin{definition}[meters Per Second Value]
  \label{def:physics:phyx-mini-0173:phyxminiproblems-problemphyxmini0173-meterspersecondvalue}
  \lean{PhyXMiniProblems.ProblemPhyXMini0173.metersPerSecondValue}
  Read a Physlib speed in metres per second
\end{definition}
\begin{proof}
  This is the declared model definition of meters Per Second Value.
\end{proof}
\begin{definition}[Figure Point]
  \label{def:physics:phyx-mini-0173:phyxminiproblems-problemphyxmini0173-figurepoint}
  \lean{PhyXMiniProblems.ProblemPhyXMini0173.FigurePoint}
  The three point labels printed in the primary figure
\end{definition}
\begin{proof}
  This is the declared model definition of Figure Point.
\end{proof}
\begin{definition}[Speaker Label]
  \label{def:physics:phyx-mini-0173:phyxminiproblems-problemphyxmini0173-speakerlabel}
  \lean{PhyXMiniProblems.ProblemPhyXMini0173.SpeakerLabel}
  The two loudspeakers displayed at A and B
\end{definition}
\begin{proof}
  This is the declared model definition of Speaker Label.
\end{proof}
\begin{definition}[figure Point]
  \label{def:physics:phyx-mini-0173:phyxminiproblems-problemphyxmini0173-speakerlabel-figurepoint}
  \lean{PhyXMiniProblems.ProblemPhyXMini0173.SpeakerLabel.figurePoint}
  \uses{def:physics:phyx-mini-0173:phyxminiproblems-problemphyxmini0173-figurepoint, def:physics:phyx-mini-0173:phyxminiproblems-problemphyxmini0173-speakerlabel}
  The figure point occupied by each loudspeaker
\end{definition}
\begin{proof}
  This is the declared model definition of figure Point.
\end{proof}
\begin{definition}[Amplifier Label]
  \label{def:physics:phyx-mini-0173:phyxminiproblems-problemphyxmini0173-amplifierlabel}
  \lean{PhyXMiniProblems.ProblemPhyXMini0173.AmplifierLabel}
  Amplifier identities used to state that both speakers share one drive
\end{definition}
\begin{proof}
  This is the declared model definition of Amplifier Label.
\end{proof}
\begin{definition}[Answer Choice]
  \label{def:physics:phyx-mini-0173:phyxminiproblems-problemphyxmini0173-answerchoice}
  \lean{PhyXMiniProblems.ProblemPhyXMini0173.AnswerChoice}
  Labels of the four answers printed with the problem
\end{definition}
\begin{proof}
  This is the declared model definition of Answer Choice.
\end{proof}
\begin{definition}[answer Frequency Hertz]
  \label{def:physics:phyx-mini-0173:phyxminiproblems-problemphyxmini0173-answerfrequencyhertz}
  \lean{PhyXMiniProblems.ProblemPhyXMini0173.answerFrequencyHertz}
  \uses{def:physics:phyx-mini-0173:phyxminiproblems-problemphyxmini0173-answerchoice}
  The frequency readout, in hertz, printed beside each answer label
\end{definition}
\begin{proof}
  This is the declared model definition of answer Frequency Hertz.
\end{proof}
\begin{definition}[recorded Dataset Answer]
  \label{def:physics:phyx-mini-0173:phyxminiproblems-problemphyxmini0173-recordeddatasetanswer}
  \lean{PhyXMiniProblems.ProblemPhyXMini0173.recordedDatasetAnswer}
  \uses{def:physics:phyx-mini-0173:phyxminiproblems-problemphyxmini0173-answerchoice}
  The answer label recorded by the supplied dataset
\end{definition}
\begin{proof}
  This is the declared model definition of recorded Dataset Answer.
\end{proof}
\begin{definition}[Two Speaker Interference Setup]
  \label{def:physics:phyx-mini-0173:phyxminiproblems-problemphyxmini0173-twospeakerinterferencesetup}
  \lean{PhyXMiniProblems.ProblemPhyXMini0173.TwoSpeakerInterferenceSetup}
  \uses{def:physics:phyx-mini-0173:phyxminiproblems-problemphyxmini0173-lengthquantity, def:physics:phyx-mini-0173:phyxminiproblems-problemphyxmini0173-frequencyquantity, def:physics:phyx-mini-0173:phyxminiproblems-problemphyxmini0173-figurepoint, def:physics:phyx-mini-0173:phyxminiproblems-problemphyxmini0173-speakerlabel, def:physics:phyx-mini-0173:phyxminiproblems-problemphyxmini0173-amplifierlabel, def:physics:phyx-mini-0173:phyxminiproblems-problemphyxmini0173-answerchoice}
  The complete apparatus and the quantities needed by the two-source model. frequencyResponse is a dimensionless response readout and lets the premise state that the two physical loudspeakers are identical without identifying a loudspeaker with a scalar. The microphone observable depends on both the chosen frequency and its positive displacement from B
\end{definition}
\begin{proof}
  This is the declared model definition of Two Speaker Interference Setup.
\end{proof}
\begin{definition}[Coherently Driven Identical Speakers]
  \label{def:physics:phyx-mini-0173:phyxminiproblems-problemphyxmini0173-coherentlydrivenidenticalspeakers}
  \lean{PhyXMiniProblems.ProblemPhyXMini0173.CoherentlyDrivenIdenticalSpeakers}
  \uses{def:physics:phyx-mini-0173:phyxminiproblems-problemphyxmini0173-twospeakerinterferencesetup}
  The speakers have the same frequency response, share an amplifier, and emit at the stated common frequency with equal phase
\end{definition}
\begin{proof}
  This is the declared model definition of Coherently Driven Identical Speakers.
\end{proof}
\begin{definition}[Has Depicted Geometry]
  \label{def:physics:phyx-mini-0173:phyxminiproblems-problemphyxmini0173-hasdepictedgeometry}
  \lean{PhyXMiniProblems.ProblemPhyXMini0173.HasDepictedGeometry}
  \uses{def:physics:phyx-mini-0173:phyxminiproblems-problemphyxmini0173-lengthquantity, def:physics:phyx-mini-0173:phyxminiproblems-problemphyxmini0173-metersvalue, def:physics:phyx-mini-0173:phyxminiproblems-problemphyxmini0173-twospeakerinterferencesetup}
  The coordinate chart is measured in metres. The clauses put B at the origin, A on the positive vertical axis, and C on the positive horizontal axis. Hence ray BC is perpendicular to segment AB. A nonnegative physical displacement x places the microphone at coordinate (x, 0)
\end{definition}
\begin{proof}
  This is the declared model definition of Has Depicted Geometry.
\end{proof}
\begin{definition}[Has Stated Readouts]
  \label{def:physics:phyx-mini-0173:phyxminiproblems-problemphyxmini0173-hasstatedreadouts}
  \lean{PhyXMiniProblems.ProblemPhyXMini0173.HasStatedReadouts}
  \uses{def:physics:phyx-mini-0173:phyxminiproblems-problemphyxmini0173-metersvalue, def:physics:phyx-mini-0173:phyxminiproblems-problemphyxmini0173-hertzvalue, def:physics:phyx-mini-0173:phyxminiproblems-problemphyxmini0173-meterspersecondvalue, def:physics:phyx-mini-0173:phyxminiproblems-problemphyxmini0173-answerchoice, def:physics:phyx-mini-0173:phyxminiproblems-problemphyxmini0173-answerfrequencyhertz, def:physics:phyx-mini-0173:phyxminiproblems-problemphyxmini0173-twospeakerinterferencesetup}
  The numerical text and answer-table readouts: separation 2.00 m, reference frequency 784 Hz, speed of sound 344 m/s, and the four displayed choices. No cutoff or interference outcome occurs in this predicate
\end{definition}
\begin{proof}
  This is the declared model definition of Has Stated Readouts.
\end{proof}
\begin{definition}[Has Physical Parameters]
  \label{def:physics:phyx-mini-0173:phyxminiproblems-problemphyxmini0173-hasphysicalparameters}
  \lean{PhyXMiniProblems.ProblemPhyXMini0173.HasPhysicalParameters}
  \uses{def:physics:phyx-mini-0173:phyxminiproblems-problemphyxmini0173-lengthquantity, def:physics:phyx-mini-0173:phyxminiproblems-problemphyxmini0173-frequencyquantity, def:physics:phyx-mini-0173:phyxminiproblems-problemphyxmini0173-metersvalue, def:physics:phyx-mini-0173:phyxminiproblems-problemphyxmini0173-hertzvalue, def:physics:phyx-mini-0173:phyxminiproblems-problemphyxmini0173-meterspersecondvalue, def:physics:phyx-mini-0173:phyxminiproblems-problemphyxmini0173-speakerlabel, def:physics:phyx-mini-0173:phyxminiproblems-problemphyxmini0173-answerchoice, def:physics:phyx-mini-0173:phyxminiproblems-problemphyxmini0173-twospeakerinterferencesetup}
  Positivity and nonnegativity conditions selecting the physical branch
\end{definition}
\begin{proof}
  This is the declared model definition of Has Physical Parameters.
\end{proof}
\begin{definition}[path Difference Meters]
  \label{def:physics:phyx-mini-0173:phyxminiproblems-problemphyxmini0173-pathdifferencemeters}
  \lean{PhyXMiniProblems.ProblemPhyXMini0173.pathDifferenceMeters}
  \uses{def:physics:phyx-mini-0173:phyxminiproblems-problemphyxmini0173-lengthquantity, def:physics:phyx-mini-0173:phyxminiproblems-problemphyxmini0173-metersvalue, def:physics:phyx-mini-0173:phyxminiproblems-problemphyxmini0173-twospeakerinterferencesetup}
  The magnitude of the two source-to-microphone path difference in metres
\end{definition}
\begin{proof}
  This is the declared model definition of path Difference Meters.
\end{proof}
\begin{definition}[Satisfies Two Source Acoustic Laws]
  \label{def:physics:phyx-mini-0173:phyxminiproblems-problemphyxmini0173-satisfiestwosourceacousticlaws}
  \lean{PhyXMiniProblems.ProblemPhyXMini0173.SatisfiesTwoSourceAcousticLaws}
  \uses{def:physics:phyx-mini-0173:phyxminiproblems-problemphyxmini0173-lengthquantity, def:physics:phyx-mini-0173:phyxminiproblems-problemphyxmini0173-frequencyquantity, def:physics:phyx-mini-0173:phyxminiproblems-problemphyxmini0173-metersvalue, def:physics:phyx-mini-0173:phyxminiproblems-problemphyxmini0173-hertzvalue, def:physics:phyx-mini-0173:phyxminiproblems-problemphyxmini0173-meterspersecondvalue, def:physics:phyx-mini-0173:phyxminiproblems-problemphyxmini0173-speakerlabel, def:physics:phyx-mini-0173:phyxminiproblems-problemphyxmini0173-twospeakerinterferencesetup, def:physics:phyx-mini-0173:phyxminiproblems-problemphyxmini0173-coherentlydrivenidenticalspeakers, def:physics:phyx-mini-0173:phyxminiproblems-problemphyxmini0173-pathdifferencemeters}
  The governing acoustic laws, uniformly stated for every positive frequency and microphone displacement: c = f λ in SI readouts; each ray length is the Euclidean source-to-microphone distance; equal-phase sources cancel exactly at odd half-wavelength path differences. These laws do not identify any answer choice as a cutoff
\end{definition}
\begin{proof}
  This is the declared model definition of Satisfies Two Source Acoustic Laws.
\end{proof}
\begin{definition}[Has No Destructive Position On BC]
  \label{def:physics:phyx-mini-0173:phyxminiproblems-problemphyxmini0173-hasnodestructivepositiononbc}
  \lean{PhyXMiniProblems.ProblemPhyXMini0173.HasNoDestructivePositionOnBC}
  \uses{def:physics:phyx-mini-0173:phyxminiproblems-problemphyxmini0173-lengthquantity, def:physics:phyx-mini-0173:phyxminiproblems-problemphyxmini0173-frequencyquantity, def:physics:phyx-mini-0173:phyxminiproblems-problemphyxmini0173-metersvalue, def:physics:phyx-mini-0173:phyxminiproblems-problemphyxmini0173-twospeakerinterferencesetup}
  There is no destructive-interference point on the open ray beyond B
\end{definition}
\begin{proof}
  This is the declared model definition of Has No Destructive Position On BC.
\end{proof}
\begin{definition}[Is No Destructive Interference Cutoff]
  \label{def:physics:phyx-mini-0173:phyxminiproblems-problemphyxmini0173-isnodestructiveinterferencecutoff}
  \lean{PhyXMiniProblems.ProblemPhyXMini0173.IsNoDestructiveInterferenceCutoff}
  \uses{def:physics:phyx-mini-0173:phyxminiproblems-problemphyxmini0173-lengthquantity, def:physics:phyx-mini-0173:phyxminiproblems-problemphyxmini0173-frequencyquantity, def:physics:phyx-mini-0173:phyxminiproblems-problemphyxmini0173-metersvalue, def:physics:phyx-mini-0173:phyxminiproblems-problemphyxmini0173-hertzvalue, def:physics:phyx-mini-0173:phyxminiproblems-problemphyxmini0173-twospeakerinterferencesetup, def:physics:phyx-mini-0173:phyxminiproblems-problemphyxmini0173-hasnodestructivepositiononbc}
  cutoff is the exact upper frequency having no destructive point: every positive frequency at or below it is safe, while every larger frequency has at least one destructive position on the open ray. Excluding displacement zero reflects the wording that the microphone is moved *out from* B and makes the boundary frequency itself part of the no-interference regime
\end{definition}
\begin{proof}
  This is the declared model definition of Is No Destructive Interference Cutoff.
\end{proof}
\begin{lemma}[wavelength meters eq speed div frequency]
  \label{lem:physics:phyx-mini-0173:phyxminiproblems-problemphyxmini0173-wavelength-meters-eq-speed-div-frequency}
  \lean{PhyXMiniProblems.ProblemPhyXMini0173.wavelength_meters_eq_speed_div_frequency}
  \uses{def:physics:phyx-mini-0173:phyxminiproblems-problemphyxmini0173-frequencyquantity, def:physics:phyx-mini-0173:phyxminiproblems-problemphyxmini0173-metersvalue, def:physics:phyx-mini-0173:phyxminiproblems-problemphyxmini0173-hertzvalue, def:physics:phyx-mini-0173:phyxminiproblems-problemphyxmini0173-meterspersecondvalue, def:physics:phyx-mini-0173:phyxminiproblems-problemphyxmini0173-twospeakerinterferencesetup, def:physics:phyx-mini-0173:phyxminiproblems-problemphyxmini0173-satisfiestwosourceacousticlaws}
  The dispersion relation determines the wavelength at any positive frequency
\end{lemma}
\begin{proof}
  The conclusion follows from the governing relations and figure readouts named in the statement of wavelength meters eq speed div frequency.
\end{proof}
\begin{lemma}[path Difference Meters eq]
  \label{lem:physics:phyx-mini-0173:phyxminiproblems-problemphyxmini0173-pathdifferencemeters-eq}
  \lean{PhyXMiniProblems.ProblemPhyXMini0173.pathDifferenceMeters_eq}
  \uses{def:physics:phyx-mini-0173:phyxminiproblems-problemphyxmini0173-lengthquantity, def:physics:phyx-mini-0173:phyxminiproblems-problemphyxmini0173-metersvalue, def:physics:phyx-mini-0173:phyxminiproblems-problemphyxmini0173-twospeakerinterferencesetup, def:physics:phyx-mini-0173:phyxminiproblems-problemphyxmini0173-hasdepictedgeometry, def:physics:phyx-mini-0173:phyxminiproblems-problemphyxmini0173-hasphysicalparameters, def:physics:phyx-mini-0173:phyxminiproblems-problemphyxmini0173-pathdifferencemeters, def:physics:phyx-mini-0173:phyxminiproblems-problemphyxmini0173-satisfiestwosourceacousticlaws}
  At a microphone displacement x, the upper path is sqrt(separation² + x²) and the lower path is x
\end{lemma}
\begin{proof}
  The conclusion follows from the governing relations and figure readouts named in the statement of path Difference Meters eq.
\end{proof}
\begin{lemma}[geometric cutoff hertz eq answer A]
  \label{lem:physics:phyx-mini-0173:phyxminiproblems-problemphyxmini0173-geometric-cutoff-hertz-eq-answera}
  \lean{PhyXMiniProblems.ProblemPhyXMini0173.geometric_cutoff_hertz_eq_answerA}
  \uses{def:physics:phyx-mini-0173:phyxminiproblems-problemphyxmini0173-metersvalue, def:physics:phyx-mini-0173:phyxminiproblems-problemphyxmini0173-meterspersecondvalue, def:physics:phyx-mini-0173:phyxminiproblems-problemphyxmini0173-answerfrequencyhertz, def:physics:phyx-mini-0173:phyxminiproblems-problemphyxmini0173-twospeakerinterferencesetup, def:physics:phyx-mini-0173:phyxminiproblems-problemphyxmini0173-hasstatedreadouts}
  The general geometric boundary c / (2 d) evaluates to 86 Hz for the stated sound speed and speaker separation. This is derived from data and is not a premise field
\end{lemma}
\begin{proof}
  The conclusion follows from the governing relations and figure readouts named in the statement of geometric cutoff hertz eq answer A.
\end{proof}
% --- Archon declaration topology end ---
% --- Archon physics formalization source end ---
